% Actual low coherent cohomology of the specified norm fourfold.
% GPT-6 Astra Pro; /root/fixed_surface; 22 November 2026 UTC.
\section[HD-F1-S21: irregularity eight and the excess boundary map]{HD-F1-S21: the actual norm fourfold has irregularity eight, with an exact excess map for its remaining scalar degree-two classes}
\label{sec:f1-norm-fourfold-low-cohomology}

For the actual smooth component $T$ in S13--S20, with
$b:T\to Y\subset P$ and $E=\mathcal O_Y$, one has
\[
 h^0(T,\mathcal O_T)=1,\qquad h^1(T,\mathcal O_T)=8.
 \tag{F1S21.1}
\]
The actual restriction map
\[
 H^1(P,\mathcal O_P)\longrightarrow H^1(T,\mathcal O_T)
                  =H^1(Y,\mathcal O_Y)
 \tag{F1S21.2}
\]
is an isomorphism. The homomorphism of Albanese varieties
induced by $T\to P$ is an isomorphism. In degree two we
obtain the explicit connecting map in (F1S21.26) below;
its kernel is exactly the part of scalar cohomology not
yet determined by the ambient twenty-eight classes.
The cohomology of its source is not a substitute for
computing that kernel.

\paragraph{1. The precise auxiliary morphism and its three flat lines.}
Retain the connected etale cyclic cover $\pi:Z\to C$ of
degree three, with $g(C)=5$ and $g(Z)=13$. Put
\[
 X=\operatorname{Sym}^8 Z,\qquad A=\operatorname{Pic}^8 C,
 \qquad F(D)=\mathcal O_C(\pi_*D).
 \tag{F1S21.3}
\]
A choice of origin identifies the torsor $A$ with the
Jacobian of $C$. Choose an Abel embedding
$a_C:C\to A$, $p\mapsto\mathcal O_C(p+D_7)$, for a fixed
degree-seven divisor $D_7$. The induced isomorphism
$a_C^*:\operatorname{Pic}^0(A)\to\operatorname{Pic}^0(C)$
specifies all identifications of flat lines below. It is
independent of the translation choice at the level of
Picard varieties.

Write $\eta_C$ for the nontrivial order-three line defining
the cover, so that
$\pi_*\mathcal O_Z=\mathcal O_C\oplus\eta_C\oplus\eta_C^2$,
after choosing one of the two character conventions. Set
\[
 \Sigma=(a_C^*)^{-1}\{\mathcal O_C,\eta_C,\eta_C^2\}
                  \subset\operatorname{Pic}^0(A).
 \tag{F1S21.4}
\]
These are exactly the $\alpha$ for which
$\pi^*a_C^*\alpha$ is trivial. Indeed projection formula
expresses its space of sections as the direct sum of the
spaces of sections of the three degree-zero lines
$a_C^*\alpha\otimes\eta_C^j$. A degree-zero line on a smooth
connected curve has a section exactly when it is trivial.
This proves the assertion, including exclusion of any
additional kernel points.

The map $F$ is surjective and has generic relative dimension
three. The finite surjection
$\operatorname{Sym}^8 Z\to\operatorname{Sym}^8 C$ factors it
through the usual Abel map. Riemann--Roch on $C$ gives
$h^0(C,L)=4$ for every $L\in\operatorname{Pic}^8 C$ except
$L=K_C$, where it is five. Thus its generic Abel fibre is
$\mathbb P^3$, while the particular fibre used here is
\[
 \mathcal T=F^{-1}(K_C)=T_0\amalg T_1\amalg T_2.
 \tag{F1S21.5}
\]
S13 proves this equality with the scheme structure, the
smoothness of each four-dimensional component, and the
transitive permutation of the components by the actual
deck transformation of $Z$. The latter also follows from
the sum-of-sheets labels in its monodromy proof: translating
all eight sheets changes that label by $8=2$ in
$\mathbb Z/3$. In particular their coherent cohomology
dimensions agree in every degree.

\paragraph{2. Twisted cohomology on the symmetric product, with its actual descent.}
For $\alpha\in\operatorname{Pic}^0(A)$ put
$L=\pi^*a_C^*\alpha$. The addition law for the normalized
Poincare line gives, for the quotient $q:Z^8\to X$,
\[
 q^*F^*\alpha\simeq L\boxtimes\cdots\boxtimes L.
 \tag{F1S21.6}
\]
Its $S_8$-linearization is the permutation of the eight
factors. This specifies the descended line, rather than
only its numerical class. Since $q$ is finite and invariants
are exact in characteristic zero,
\[
 R\Gamma(X,F^*\alpha)\simeq
       \bigl(R\Gamma(Z,L)^{\otimes8}\bigr)^{S_8}.
 \tag{F1S21.7}
\]
If $\alpha\notin\Sigma$, the nontrivial degree-zero line
$L$ has $h^0(L)=0$ and $h^1(L)=12$. The Koszul signs for
the permutation of the degree-one factors therefore give
\[
 H^j(X,F^*\alpha)=0\ (j\ne8),\qquad
 H^8(X,F^*\alpha)=\bigwedge^8 H^1(Z,L).
 \tag{F1S21.8}
\]
If $\alpha\in\Sigma$, a trivialization of $L$ gives an
$S_8$-invariant trivialization of its eightfold external
product, so $F^*\alpha\simeq\mathcal O_X$. The same
super-symmetric calculation then gives
\[
 h^j(X,\mathcal O_X)=\binom{13}{j}\quad(0\le j\le8).
 \tag{F1S21.9}
\]
Serre duality in dimension eight converts (F1S21.8) into
\[
 H^j(X,\omega_X\otimes F^*\alpha)=0
       \quad(j>0,\ \alpha\notin\Sigma).
 \tag{F1S21.10}
\]

\paragraph{3. All positive higher direct images are sums of the three specified lines.}
Let $\mathcal F_i=R^iF_*\omega_X$. The torsion-freeness
and splitting theorems for canonical direct images give
\[
 \mathcal F_i=0\ (i>3),\qquad
 RF_*\omega_X\simeq\bigoplus_{i=0}^3\mathcal F_i[-i].
 \tag{F1S21.11}
\]
Every hypothesis is present: $X$ is smooth projective,
$F$ is proper and surjective, $A$ is smooth, and the
generic relative dimension is three. In particular
torsion freeness excludes a higher direct image supported
only at the exceptional four-dimensional fibre.

For $i>0$, the projection formula and (F1S21.11) make
$H^0(A,\mathcal F_i\otimes\alpha)$ a direct summand of
$H^i(X,\omega_X\otimes F^*\alpha)$. Thus (F1S21.10)
implies
\[
 V^0(\mathcal F_i):=
 \{\alpha:H^0(A,\mathcal F_i\otimes\alpha)\ne0\}
                           \subseteq\Sigma\quad(i>0).
 \tag{F1S21.12}
\]
The Chen--Jiang decomposition for \emph{each} higher
canonical direct image writes $\mathcal F_i$ as a direct
sum of $\gamma\otimes q^*\mathcal G$, where $q:A\to B$
is a quotient with connected fibres, $\gamma$ is a
torsion flat line, and $\mathcal G$ is a nonzero
M-regular coherent sheaf on $B$.

Here the finite support in (F1S21.12) forces every such
$B$ to be a point. To see the exact implication, a nonzero
M-regular sheaf has $V^0(\mathcal G)=\operatorname{Pic}^0(B)$:
the GV support-locus property says $V^0$ is nonempty and
every component of codimension $k$ is contained in $V^k$;
M-regularity excludes $k\ge1$. As $q_*\mathcal O_A=\mathcal O_B$,
projection formula identifies
\[
 H^0(A,q^*(\mathcal G\otimes\delta))
                      =H^0(B,\mathcal G\otimes\delta).
 \tag{F1S21.13}
\]
Consequently the support locus of $\gamma\otimes q^*\mathcal G$
contains $\gamma^{-1}q^*\operatorname{Pic}^0(B)$.
The dual of the connected-fibre quotient is injective,
so a positive-dimensional $B$ contradicts (F1S21.12).
For a point $B$, the summand is a sum of copies of
$\gamma$, whose $V^0$ is the single point $\gamma^{-1}$.
It follows that
\[
 \mathcal F_i\text{ is a direct sum of lines in }\Sigma,
                         \qquad i=1,2,3.
 \tag{F1S21.14}
\]
No local freeness assertion for $\mathcal F_0$ is used.

The precise primary inputs are Pareschi--Popa--Schnell,
\emph{Hodge modules on complex tori and generic vanishing
for compact K\"ahler manifolds}, Geometry \& Topology
21 (2017), 2419--2460, Theorem A (p.~2420),
Proposition~13.6 (p.~2446), Theorems~14.1--14.2
(pp.~2446--2447), and the proof of Theorem A in
Section~16 (p.~2449). Theorem A covers all $i\ge0$;
its proof passes from the canonical direct-image Hodge
module to Theorem D. These statements and the complete
relevant Section~15--16 proofs were inspected in the
author and publisher texts.

\paragraph{4. Duality fixes the shifts and the multiplicities.}
Write $\mathbb D_A(-)=R\mathcal Hom_A(-,\mathcal O_A)$.
Grothendieck duality, with $\dim X=8$, $\dim A=5$ and
$\omega_A\simeq\mathcal O_A$, gives the exact isomorphism
\[
 \mathbb D_A(RF_*\mathcal O_X)\simeq RF_*\omega_X[3].
 \tag{F1S21.15}
\]
Bidualizing (F1S21.11) therefore yields
\[
 RF_*\mathcal O_X\simeq
 \mathbb D_A(\mathcal F_0)[-3]
 \oplus\mathcal F_1^\vee[-2]
 \oplus\mathcal F_2^\vee[-1]
 \oplus\mathcal F_3^\vee.
 \tag{F1S21.16}
\]
The first summand contributes nothing in degrees below
three, both to global hypercohomology and to a derived
fibre at any point of $A$. For the latter assertion,
take a finite locally free resolution of $\mathcal F_0$
in nonpositive degrees. Its dual lies in nonnegative
degrees and stays there after tensoring with a residue
field. The shift $[-3]$ places it in degrees at least
three. This argument does not replace derived tensor
with ordinary tensor for a possibly nonflat sheaf.

For $\alpha\in\Sigma$, let $m_i(\alpha)$ be the multiplicity
of $\alpha^{-1}$ in $\mathcal F_i^\vee$. A nontrivial flat
line on an abelian variety has zero coherent cohomology
in all degrees; the trivial line on $A$ has cohomology
dimensions $\binom5j$. Tensor (F1S21.16) by $\alpha$ and
use (F1S21.9). The successive degree-zero, degree-one and
degree-two equalities are
\[
 \begin{split}
 1&=m_3(\alpha),\\
 13&=5m_3(\alpha)+m_2(\alpha),\\
 78&=10m_3(\alpha)+5m_2(\alpha)+m_1(\alpha).
 \end{split}
 \tag{F1S21.17}
\]
Thus, for \emph{each} of the three points $\alpha$,
the respective multiplicities are $1,8,28$. In particular
\[
 \mathcal F_3^\vee\simeq\bigoplus_{\alpha\in\Sigma}\alpha,
 \quad
 \mathcal F_2^\vee\simeq\bigoplus_{\alpha\in\Sigma}\alpha^{\oplus8},
 \quad
 \mathcal F_1^\vee\simeq\bigoplus_{\alpha\in\Sigma}\alpha^{\oplus28}.
 \tag{F1S21.18}
\]
Taking the derived fibre of (F1S21.16) at $K_C$ therefore
gives dimensions $3,24,84$ in degrees $0,1,2$.

\paragraph{5. The actual derived fibre and its excess line.}
Let $k_{K_C}$ denote the skyscraper sheaf of the specified
point in $A$, and put
\[
                      \mathcal K=LF^*k_{K_C}\in D^b(X).
 \tag{F1S21.19}
\]
This object is computed near $\mathcal T$ by the Koszul
complex on the five pulled-back local parameters of
$A$. S13 gives the actual conormal sequence
\[
 0\longrightarrow\mathcal O_{\mathcal T}(-h)
 \longrightarrow T^*_{K_C}A\otimes\mathcal O_{\mathcal T}
 \longrightarrow N^*_{\mathcal T/X}\longrightarrow0,
 \qquad h=p^*\mathcal O_{\mathbb P^4}(1)=3\eta.
 \tag{F1S21.20}
\]
Here $p:\mathcal T\to|K_C|$ is the actual finite morphism,
and the first line is generated at $D$ by the canonical
section with zero divisor $\pi_*D$. At a point of
$\mathcal T$, four of the five equations generate its
regular-embedding ideal. Invertible row operations reduce
the Koszul sequence to four regular generators and a
zero fifth generator. Their homology, glued by the kernel
line in (F1S21.20), gives
\[
 \mathcal H^0(\mathcal K)=\mathcal O_{\mathcal T},\qquad
 \mathcal H^{-1}(\mathcal K)=\mathcal O_{\mathcal T}(-h),
 \qquad \mathcal H^r(\mathcal K)=0\quad(r\ne0,-1).
 \tag{F1S21.21}
\]
This determines the homology sheaves; a splitting of
$\mathcal K$ into them is not required.

The line $h$ is ample because $p$ is finite. On each smooth
fourfold component, Kodaira vanishing and Serre duality
give
\[
 H^q(T_j,\mathcal O_{T_j}(-h))
 \simeq H^{4-q}(T_j,\omega_{T_j}(h))^*=0
                                  \qquad(q<4).
 \tag{F1S21.22}
\]
The canonical truncation triangle
\[
 \mathcal O_{\mathcal T}(-h)[1]\longrightarrow\mathcal K
 \longrightarrow\mathcal O_{\mathcal T}
 \longrightarrow\mathcal O_{\mathcal T}(-h)[2]
 \tag{F1S21.23}
\]
therefore induces actual isomorphisms
\[
 H^j(R\Gamma(X,\mathcal K))
       \xrightarrow{\ \sim\ }H^j(\mathcal T,\mathcal O_{\mathcal T})
                           \qquad(j=0,1).
 \tag{F1S21.24}
\]

There is also a direct comparison with the derived fibre
of (F1S21.16). Since $k_{K_C}$ is perfect on the smooth
variety $A$, the derived projection formula gives
\[
 \begin{split}
 R\Gamma(X,\mathcal K)
  &\simeq R\Gamma\bigl(A,
             RF_*\mathcal O_X\otimes^L k_{K_C}\bigr)\\
  &\simeq L\iota_{K_C}^*(RF_*\mathcal O_X).
 \end{split}
 \tag{F1S21.25}
\]
In the second line both sides are complexes of vector
spaces; the sheaf in the first line is the pushforward
from the specified point of the displayed derived fibre.
This comparison uses the actual derived pullback, so it
does not require $F$ to be flat at $K_C$ or assert a
Tor-independent ordinary fibre square.

Equations (F1S21.18), (F1S21.24) and (F1S21.25) prove
$h^0(\mathcal O_{\mathcal T})=3$ and
$h^1(\mathcal O_{\mathcal T})=24$. The actual deck
permutation identifies the three component cohomology
spaces, so division by three proves (F1S21.1).
In the next degree, the \emph{same} triangle gives the
precise exact sequence
\[
 0\longrightarrow H^2(R\Gamma(X,\mathcal K))
 \longrightarrow H^2(\mathcal T,\mathcal O_{\mathcal T})
 \longrightarrow H^4(\mathcal T,\mathcal O_{\mathcal T}(-h))
 \xrightarrow{\delta}
 H^3(R\Gamma(X,\mathcal K)),
 \tag{F1S21.26}
\]
where the left vector space has dimension $84$. The map
$\delta$ is induced by the first arrow of the actual
truncation triangle (F1S21.23). The right-hand target
identifies with
\[
 H^3\bigl(L\iota_{K_C}^*RF_*\mathcal O_X\bigr)
   =H^0\bigl(L\iota_{K_C}^*\mathbb D_A(\mathcal F_0)\bigr)
   =\operatorname{Hom}_A(\mathcal F_0,k_{K_C}).
 \tag{F1S21.26a}
\]
For the last equality, dualize a finite locally free
resolution of $\mathcal F_0$ and tensor with the residue
field, as in paragraph~4.

In particular the group $H^4(\mathcal O_{\mathcal T}(-h))$
cannot be dropped from (F1S21.26). It occurs on the right
of the degree-two comparison despite Kodaira vanishing
in degrees below four. A provisional draft claiming a
degree-two isomorphism from (F1S21.22) was withdrawn
before promotion; (F1S21.26) is the corrected statement.

\paragraph{6. The actual Albanese map and the remaining scalar cohomology.}
S18 proves that restriction from $P$ has rank eight in
degree one. The dimension just computed makes this very
map an isomorphism. S16--S17 identify $Rb_*\mathcal O_T$
with $\mathcal O_Y$, proving (F1S21.2) with its asserted
target.

The homomorphism $\operatorname{Alb}(T)\to P$ induced by
the given Abel map is surjective with connected kernel
by S19. Its source now has dimension $q(T)=8=\dim P$.
Its connected kernel has dimension zero and is therefore
trivial, proving the asserted isomorphism of the actual
abelian varieties, not only of their tangent spaces.

The image of the left arrow in (F1S21.26) is exactly the
direct sum of the three ambient degree-two subspaces.
Here is the multiplicative identification. The object
$\mathcal K=\mathcal O_X\otimes^L_{\mathcal O_A}k_{K_C}$
is a derived commutative algebra, and its augmentation
to $\mathcal O_{\mathcal T}$ preserves multiplication.
The degree-zero isomorphism identifies its three
idempotents with the indicators of $T_0,T_1,T_2$.
Its degree-one isomorphism lifts every element of
$H^1(\mathcal O_{T_j})=H^1(\mathcal O_P)$ to the
corresponding idempotent part of $H^1(R\Gamma\mathcal K)$.
Taking their products shows that the image in degree
two contains the ambient subspace
$\bigwedge^2H^1(P,\mathcal O_P)$ on every component.
S18 proves that each of these three subspaces has
dimension twenty-eight. Their direct sum has dimension
$84$, which is already the full dimension of the
injecting space in (F1S21.26). This proves the claimed
identification through the actual algebra map.
Consequently
\[
 \frac{H^2(\mathcal T,\mathcal O_{\mathcal T})}
 {\displaystyle\bigoplus_{j=0}^2
       \operatorname{im}\bigl(H^2(P,\mathcal O_P)
                    \to H^2(T_j,\mathcal O_{T_j})\bigr)}
                          \simeq\ker\delta.
 \tag{F1S21.27}
\]
The deck transformation preserves the triangle and
permutes its component parts. It follows also that
\[
                 h^2(T,\mathcal O_T)
                       =28+\frac{\dim\ker\delta}{3}.
 \tag{F1S21.28}
\]
Equation (F1S21.28) is an exact computation in terms of
the displayed map for this actual fibre, not a
conditional semiregularity assertion. We have not
proved that $\delta$ is injective.

S18's unit/multiplication retract still identifies
$H^2(Y,\mathcal O_Y)$ with a specified direct summand
of $\operatorname{Ext}^2_P(E,E)$, and the ambient
$28$-dimensional part has injective full trace. The
extra scalar quotient (F1S21.27), together with the
other self-Ext directions, remains in the full
trace calculation. Neither the derived decomposition
(F1S21.16) nor the irregularity computation discards it.
